\section{Related Work}
\label{sec:related}

%TL;DR Trained block-world generators learn structure from a corpus; we test a training-free agent against them and find the trained model still ahead on geometry.
\paragraph{Trained generators for block worlds.}
Learned generators for Minecraft write a building block by block in the order a player placed
it~\citep{chen2019voxelcnn}, denoise whole chunks after training on billions of
cubes~\citep{merino2026dreamcubed}, or optimize a block field against a text-image score for each
prompt~\citep{earle2024dreamcraft}. For bricks, \citet{pun2025brickgpt} train an autoregressive
model on a text-annotated brick corpus and reject unstable placements during decoding. All of these
learn structure from data or from per-prompt optimization, and each new kind of condition needs a
new corpus or objective. Our agent needs neither, which makes it cheap to condition on an imagined
image; we keep a trained voxel transformer in the comparison to show what the agent gives up in
geometry.

%TL;DR LLM building agents in Minecraft take text, a plan or a given reference image; none measures whether a self-drawn image helps against a loop matched in turns.
\paragraph{Language-model building agents in Minecraft.}
\citet{hu2024t2bm} prompt language models to write Minecraft buildings from text, and
\citet{chen2024apt} decompose the prompt into an architectural plan and blueprint with memory and
reflection modules. CraftUtopia~\citep{fu2026craftutopia} turns a given 2D reference image into a
blueprint and splits construction across parallel workers. Word2Minecraft~\citep{huang2025word2minecraft}
generates game levels from text, Voyager~\citep{wang2023voyager} writes skill code to act in the game, and
MineAnyBuild~\citep{wei2025mineanybuild} and MC-Bench~\citep{mcbench2025} evaluate agents and
models that build from instructions. In these systems the image, when present, is an input the user
provides. We ask a different question: whether an image the agent draws for itself improves the
build, measured against a loop that spends the same turns without one.

%TL;DR Code-writing 3D agents with render feedback chase a given target image or a validator; our target is self-made, so the open question is whether it beats having no target.
\paragraph{Code-writing 3D agents with visual feedback.}
Agents that write graphics programs and repair them from renders are common outside block worlds.
SceneCraft~\citep{hu2024scenecraft} and LL3M~\citep{lu2025ll3m} write Blender code for scenes and
assets, CADCodeVerify~\citep{alrashedy2025cadcodeverify} checks CAD code against renders of its
output, and VIGA~\citep{yin2026viga} and LEGO-Anything~\citep{li2026legoanything} run a
write-render-compare loop until the program reproduces a target image. For bricks,
BrickWright~\citep{pun2026brickwright} and BrickBench~\citep{kulits2026brickbench} give coding
agents a brick assembly API with a physical validator. In inverse graphics the target is given and
the render-to-target distance is the task; in our setting the target is itself generated, so a
build that matches it perfectly is only as good as the picture. Benchmarks of code-written 3D
report that executable code is often not spatially correct~\citep{zheng2026voxelcodebench} and that
successful renders often hide floating or disconnected geometry~\citep{gao2026threedcodebench}. We
find the analogous gap one level up: a build that looks closer to real houses in renders is not
closer in geometry or in walkability.

%TL;DR Imagined goals help robot and web agents because there the image fixes the goal state; one exterior view underspecifies a building, which predicts our result.
\paragraph{Imagined goals for acting.}
Generating the goal before acting helps embodied agents: SuSIE~\citep{black2024susie} edits the
current camera image into a subgoal for a low-level policy, UniPi~\citep{du2023unipi} plans as a
generated video and reads actions from it, and WebDreamer~\citep{gu2024webdreamer} uses a language
model to simulate the outcome of each web action. CaP-X~\citep{fu2026capx} studies code-writing robot
agents and lists visual differencing among the test-time techniques that narrow the gap to human-designed
abstractions. In manipulation a subgoal image specifies most of the state that matters. A single
exterior view of a house does not: it fixes the facade and leaves the interior, the back and the
circulation open. Our results match this difference: the imagined image moves appearance and leaves
structure where it was.

%TL;DR We adopt an existing evaluation protocol unchanged and make no claim about evaluation.
\paragraph{Evaluation.}
We score every arm with the protocol of a concurrent submission~\citep{anonymous2027blindness}
without modification, including its distribution metrics, its walkability probe and its novelty
gates, and we make no claim about the protocol itself.
